% Informe de sprint, semana 6 (avance del curso Proyecto de Investigación 2).
% Fuente: docs/informe_sprint_sem6.md, rama avance-sem6-sprint. Mismas cifras; cada una tiene su
% fuente en notebooks/results/sprint6_results.csv o en el documento del proyecto que se cita.
% Figuras: subir docs/figs/sprint6/fig_ab_datos.pdf y fig_anexo_modelos.pdf a E_IMAGENES/6_Sprint/.
% Paquetes usados: graphicx, float ([H]), amsmath. Las citas usan claves de docs/literature_thesis.bib.

\chapter{Informe de sprint: semana 6}
\markboth{CAPÍTULO \thechapter: Informe de sprint: semana 6}{}

Este informe reúne el análisis exploratorio de datos (semana 4), la ingeniería de características
(semana 5) y un experimento A/B de los datos de entrenamiento (semana 6). Todas las cifras provienen
de cuadernos ejecutados y de registros versionados del proyecto: el registro de corridas
(\texttt{docs/\allowbreak{}EXPERIMENT\_\allowbreak{}LOG.md}), el registro de decisiones metodológicas
(\texttt{docs/\allowbreak{}METHODS\_\allowbreak{}NOTES.md}) y la tabla de resultados con la fuente de cada valor
(\texttt{notebooks/\allowbreak{}results/\allowbreak{}sprint6\_\allowbreak{}results.csv}). La rama del entregable es
\texttt{avance-\allowbreak{}sem6-\allowbreak{}sprint}.

\begin{table}[H]
    \centering
    \small
    \caption{Material del sprint}
    \label{tab:sprint_material}
    \begin{tabular}{|p{4.2cm}|p{5.0cm}|p{5.0cm}|}
        \hline
        \textbf{Parte} & \textbf{Cuaderno (ejecutado)} & \textbf{Documento de apoyo} \\
        \hline
        Análisis exploratorio (semana 4) & \texttt{notebooks/\allowbreak{}eda.ipynb} & \texttt{docs/\allowbreak{}peru\_\allowbreak{}data\_\allowbreak{}card.md} \\
        \hline
        Ingeniería de características (semana 5) & \texttt{notebooks/\allowbreak{}feature\_\allowbreak{}engineering.ipynb} & \texttt{docs/\allowbreak{}feature\_\allowbreak{}engineering.md} \\
        \hline
        Experimento A/B de los datos (semana 6) & \texttt{notebooks/\allowbreak{}baseline\_\allowbreak{}variantes.ipynb} & \texttt{notebooks/\allowbreak{}results/\allowbreak{}sprint6\_\allowbreak{}results.csv} \\
        \hline
    \end{tabular}
\end{table}

\section{Contexto y métrica central}

El objetivo del proyecto es segmentar el agua de inundación en imágenes Sentinel-2 de Perú. El caso
realista es una inundación nueva \textbf{sin etiquetas peruanas}: solo se dispone de un conjunto
etiquetado de otras regiones (Sen1Floods11) y de imágenes de Perú sin etiquetar.

Los objetivos del sprint son tres:
\begin{enumerate}
    \item Conocer los datos y sus riesgos antes de modelar: fuga de información, desbalance de
    clases, cambio de dominio y calidad de las etiquetas (semana 4).
    \item Medir qué aportan las características diseñadas a mano frente a la reflectancia cruda
    (semana 5).
    \item Probar con un experimento A/B si \textbf{cambiar solo los datos de entrenamiento} corrige
    el problema identificado en las semanas 4 y 5 (semana 6).
\end{enumerate}

La métrica central es el IoU de agua agrupado sobre todos los píxeles etiquetados de
\texttt{peru-\allowbreak{}dev}:
\begin{equation}
    \mathrm{IoU} = \frac{TP}{TP + FP + FN}.
    \label{eq:iou_sprint}
\end{equation}
Se elige el IoU y no la exactitud porque el agua es una clase minoritaria
(sección~\ref{sec:sprint_eda}): un modelo que no detecta agua tendría una exactitud alta. Las
métricas secundarias son la precisión y el recall del agua, y el IoU en \texttt{peru-\allowbreak{}test}
(evaluación única).

\section{Datos}

\begin{table}[H]
    \centering
    \small
    \caption{Conjuntos de datos}
    \label{tab:sprint_datos}
    \begin{tabular}{|p{3.4cm}|p{3.4cm}|p{3.4cm}|p{4.0cm}|}
        \hline
        \textbf{Conjunto} & \textbf{Fuente} & \textbf{Recortes} & \textbf{Uso} \\
        \hline
        Sen1Floods11 (etiquetas manuales) & \cite{bonafilia_sen1floods11_2020} & 252 / 89 / 90 (train / valid / test oficiales) & Entrenamiento y selección del experimento A/B; no se vuelve a partir \\
        \hline
        Perú 2023, ciclón Yaku (Lambayeque y Piura) & Mapas de UNOSAT \cite{unosat_lambayequepiura22mar_2023, unosat_lambayequepiura1apr_2023} & 753 / 174 / 172 (\texttt{peru-\allowbreak{}train} / \texttt{peru-\allowbreak{}dev} / \texttt{peru-\allowbreak{}test}) & Imágenes de \texttt{peru-\allowbreak{}train} sin etiquetas para las variantes; \texttt{peru-\allowbreak{}dev} para comparar; \texttt{peru-\allowbreak{}test} una sola vez \\
        \hline
        Perú 2017, Niño Costero & Copernicus EMS EMSR199 \cite{copernicus_emsr199_2017} & 204, de ellos 101 sin contacto con \texttt{peru-\allowbreak{}train} & Prueba en otro evento \\
        \hline
    \end{tabular}
\end{table}

\begin{itemize}
    \item \textbf{Imágenes.} Sentinel-2 L1C, 13 bandas, recortes de $512 \times 512$ píxeles a 10~m
    ($5.12 \times 5.12$~km), en UTM 17S. Las imágenes de Perú se descargan de la API de Copernicus
    Data Space.
    \item \textbf{Armonización radiométrica.} Se elimina el desplazamiento de $+1000$~ND de los
    productos con \textit{processing baseline} $\geq$ 04.00, para que los valores de Perú sean
    comparables con Sen1Floods11.
    \item \textbf{Etiquetas de Perú sin anotación manual.} Los polígonos expertos se rasterizan con
    reglas fijas: 1 = agua, 0 = no agua, $-1$ = ignorado. Se ignoran los píxeles fuera del área
    analizada por UNOSAT, bajo nubes o sin datos.
    \item \textbf{Partición espacial.} Bloques de $4 \times 4$ recortes (unos 20~km), asignados al
    azar con semilla fija. Los recortes de entrenamiento vecinos de dev o prueba se descartan como
    \textit{buffer} (192 recortes). Una misma ubicación queda siempre en la misma partición en las
    dos fechas.
\end{itemize}

\section{Análisis exploratorio (semana 4)}
\label{sec:sprint_eda}

El cuaderno \texttt{notebooks/\allowbreak{}eda.ipynb} solo lee etiquetas de \texttt{peru-\allowbreak{}train},
\texttt{peru-\allowbreak{}dev} y de las particiones train y valid de Sen1Floods11. De las particiones de prueba
solo se muestran metadatos.

\subsection{Hallazgos principales}

\begin{table}[H]
    \centering
    \small
    \caption{Proporción de agua por conjunto}
    \label{tab:sprint_balance}
    \begin{tabular}{|l|c|c|}
        \hline
        \textbf{Conjunto} & \textbf{Agua (\% de píxeles etiquetados)} & \textbf{Recortes sin agua} \\
        \hline
        \texttt{peru-\allowbreak{}train} & 4.6~\% & 64.5~\% \\
        \hline
        \texttt{peru-\allowbreak{}dev} & 12.4~\% & 36.2~\% \\
        \hline
        Sen1Floods11 train & 9.5~\% & 13.1~\% \\
        \hline
        Sen1Floods11 valid & 11.0~\% & 13.5~\% \\
        \hline
    \end{tabular}
\end{table}

\begin{itemize}
    \item \textbf{Desbalance.} El agua es una clase minoritaria y su proporción cambia entre
    particiones (tabla~\ref{tab:sprint_balance}).
    \item \textbf{Fuga de información.} Ninguna ubicación aparece en dos particiones (0 casos
    encontrados).
    \item \textbf{Cambio de dominio.} El suelo seco de Perú (desierto de Sechura) es mucho más
    brillante que el de Sen1Floods11; el agua se parece más entre ambos dominios.
    \item \textbf{Calidad de etiquetas.} 16 de 753 recortes de entrenamiento y 8 de 174 de dev
    contradicen claramente al índice MNDWI.
    \item \textbf{Correlación entre bandas} (\texttt{peru-\allowbreak{}train}): igual o mayor a 0.97 dentro de cada
    grupo (visible B02--B05, borde rojo/NIR B06--B8A, SWIR B11--B12) y entre 0.70 y 0.91 entre grupos.
    \item \textbf{Nubes.} Los detectores estándar (s2cloudless, SCL) marcan el desierto como nube, así
    que se usa una regla propia de nube espesa.
\end{itemize}

\subsection{Decisiones tomadas a partir del EDA}

\begin{table}[H]
    \centering
    \small
    \caption{Decisiones derivadas del EDA}
    \label{tab:sprint_eda_decisiones}
    \begin{tabular}{|p{5.0cm}|p{9.5cm}|}
        \hline
        \textbf{Riesgo} & \textbf{Decisión} \\
        \hline
        Fuga espacial & Partición por bloques con \textit{buffer} \\
        \hline
        Desbalance & IoU de agua agrupado como métrica principal (no exactitud); precisión y recall por separado \\
        \hline
        Distinta proporción de agua entre particiones & Comparar siempre dentro de la misma partición \\
        \hline
        Cambio de dominio & Estudiar cómo transferir desde Sen1Floods11 sin etiquetas peruanas (experimento de la semana 6) \\
        \hline
        Ruido en etiquetas & Etiquetas congeladas; subconjunto sin recortes dudosos como análisis de sensibilidad \\
        \hline
    \end{tabular}
\end{table}

\section{Ingeniería de características (semana 5)}

El cuaderno \texttt{notebooks/\allowbreak{}feature\_\allowbreak{}engineering.ipynb} entrena un clasificador por píxel
(\textit{gradient boosting}) con cuatro conjuntos de variables:
\begin{itemize}
    \item las 13 bandas;
    \item 8 índices: NDWI, MNDWI \cite{xu_modification_2006}, MNDWI2, NDVI, NDMI, AWEI\_nsh, AWEI\_sh
    y WRI;
    \item bandas más índices;
    \item todo lo anterior más 5 variables de contexto local (ventanas de $3 \times 3$ y $7 \times 7$).
\end{itemize}
Se consideran dos escenarios: (a) entrenado con etiquetas de Perú y (b) entrenado solo con
Sen1Floods11. La evaluación se hace en \texttt{peru-\allowbreak{}dev}, sobre 1/4 de los píxeles, con 3 semillas.

\begin{table}[H]
    \centering
    \small
    \caption{Resultados de la ingeniería de características (IoU de agua en \texttt{peru-\allowbreak{}dev})}
    \label{tab:sprint_fe}
    \begin{tabular}{|p{3.6cm}|p{4.8cm}|p{5.8cm}|}
        \hline
        \textbf{Escenario} & \textbf{Mejor clasificador} & \textbf{Umbral de MNDWI} \\
        \hline
        (a) Etiquetas de Perú & 0.930 $\pm$ 0.001 (bandas crudas) & 0.924 (umbral 0.28, ajustado en \texttt{peru-\allowbreak{}train}) \\
        \hline
        (b) Solo Sen1Floods11 & 0.71--0.74 (según el conjunto) & 0.836 (umbral 0.15, ajustado en Sen1Floods11) \\
        \hline
    \end{tabular}
\end{table}

Conclusiones que alimentan la semana 6:
\begin{itemize}
    \item Con etiquetas del dominio, las variables diseñadas no añaden información.
    \item Sin etiquetas de Perú, los clasificadores entrenados en Sen1Floods11 confunden el desierto
    brillante con agua (precisión de 0.72 a 0.74).
    \item \textbf{El problema principal es el cambio de dominio del suelo seco, no la falta de
    variables.} Por eso el experimento de esta semana cambia los datos y no el modelo.
\end{itemize}

\section{Línea base}

La línea base, \texttt{UNet\_\allowbreak{}s1f11\_\allowbreak{}raw}, es una U-Net \cite{ronneberger_unet_2015} con codificador
ResNet-34 \cite{he_deep_2016} preentrenado en ImageNet, con 13 bandas L1C, entrenada con las
imágenes crudas de Sen1Floods11 \texttt{train}. Es lo que se obtendría hoy ante una inundación nueva
en Perú, sin etiquetas: \textbf{0.836 $\pm$ 0.022} de IoU en \texttt{peru-\allowbreak{}dev}, con precisión 0.848
y recall 0.985. Su error dominante son los falsos positivos en el desierto. En \texttt{peru-\allowbreak{}test}
obtiene 0.576.

\section{Experimentos A/B: un solo cambio por variante}

\begin{table}[H]
    \centering
    \small
    \caption{Elementos comunes a las tres filas del experimento}
    \label{tab:sprint_fijo}
    \begin{tabular}{|p{4.0cm}|p{10.5cm}|}
        \hline
        \textbf{Elemento} & \textbf{Valor fijo} \\
        \hline
        Modelo & U-Net, codificador ResNet-34 (ImageNet), 13 bandas L1C \\
        \hline
        Datos de entrenamiento & Sen1Floods11 \texttt{train} (252 recortes) \\
        \hline
        Selección & Mejor época por IoU de agua en Sen1Floods11 \texttt{valid} \\
        \hline
        Configuración & \texttt{configs/\allowbreak{}unet.yaml}: pérdida Dice, AdamW, tasa de aprendizaje $3\times10^{-4}$, lote 8, aumento D4 \\
        \hline
        Parada & Máximo 200 épocas, parada temprana con paciencia 20 \\
        \hline
        Semillas & 0, 1 y 2 \\
        \hline
        Hardware & RunPod RTX 4090, precisión mixta de 16 bits \\
        \hline
    \end{tabular}
\end{table}

Solo cambia la transformación de las imágenes de entrenamiento (y de \texttt{valid}, para que la
selección vea el mismo estilo de entrada). En los tres casos, el modelo final se aplica a las
imágenes de Perú \textbf{sin modificar}.

\begin{table}[H]
    \centering
    \small
    \caption{Variantes del experimento A/B}
    \label{tab:sprint_variantes}
    \begin{tabular}{|p{3.3cm}|p{1.4cm}|p{5.2cm}|p{4.6cm}|}
        \hline
        \textbf{ID} & \textbf{Rol} & \textbf{Qué cambia} & \textbf{Por qué} \\
        \hline
        \texttt{UNet\_\allowbreak{}s1f11\_\allowbreak{}raw} & \textit{Baseline} & Nada & Referencia sin adaptación \\
        \hline
        \texttt{UNet\_\allowbreak{}s1f11\_\allowbreak{}trHM} & Var1 & Igualación de histogramas por banda: cada valor se lleva al cuantil equivalente de las imágenes de \texttt{peru-\allowbreak{}train} sin etiquetas (1\,001 cuantiles) & Si el suelo seco peruano tiene otra distribución de brillo, igualar las distribuciones debería corregirlo \\
        \hline
        \texttt{UNet\_\allowbreak{}s1f11\_\allowbreak{}trFDA\_\allowbreak{}b0.01} & Var2 & FDA: la amplitud de Fourier de baja frecuencia ($\beta = 0.01$) de cada imagen se reemplaza por la de una imagen de \texttt{peru-\allowbreak{}train} al azar; se conserva la fase & Alternativa que transfiere solo el ``estilo'' global de la imagen y conserva su estructura \\
        \hline
    \end{tabular}
\end{table}

Las filas de referencia no son variantes, porque cambian más de una cosa: la U-Net entrenada con
etiquetas de Perú (cota superior) y el umbral de MNDWI de la semana 5.

\section{Resultados comparables}

\subsection{Tabla estándar}

\begin{table}[H]
    \centering
    \small
    \caption{Resultados del experimento A/B: métrica central y secundarias}
    \label{tab:sprint_ab_metricas}
    \begin{tabular}{|p{1.5cm}|p{3.3cm}|p{2.4cm}|p{2.6cm}|p{1.8cm}|p{2.4cm}|}
        \hline
        \textbf{Rol} & \textbf{ID} & \textbf{IoU \texttt{peru-\allowbreak{}dev} (media $\pm$ de)} & \textbf{Semillas 0 / 1 / 2} & \textbf{Precisión / recall dev} & \textbf{IoU \texttt{peru-\allowbreak{}test} [IC 95~\%]} \\
        \hline
        \textit{Baseline} & \texttt{UNet\_\allowbreak{}s1f11\_\allowbreak{}raw} & 0.8364 $\pm$ 0.0224 & 0.8433 / 0.8545 / 0.8114 & 0.848 / 0.985 & 0.576 [0.467, 0.669] \\
        \hline
        Var1 & \texttt{UNet\_\allowbreak{}s1f11\_\allowbreak{}trHM} & \textbf{0.9118 $\pm$ 0.0026} & 0.9148 / 0.9098 / 0.9108 & 0.962 / 0.946 & \textbf{0.798 [0.704, 0.865]} \\
        \hline
        Var2 & \texttt{UNet\_\allowbreak{}s1f11\_\allowbreak{}trFDA\_\allowbreak{}b0.01} & 0.6151 $\pm$ 0.1122 & 0.7251 / 0.5009 / 0.6193 & no registrado & no evaluado (descartado en dev) \\
        \hline
        Referencia & \texttt{UNet\_\allowbreak{}peru} (etiquetas de Perú) & 0.9415 $\pm$ 0.0019 & 0.9397 / 0.9434 / 0.9414 & --- & 0.828 [0.732, 0.895] \\
        \hline
        Referencia & \texttt{MNDWI\_\allowbreak{}s1f11\_\allowbreak{}t0.15} (umbral, semana 5) & 0.836 (1/4 de los píxeles) & --- & --- & --- \\
        \hline
    \end{tabular}
\end{table}

\begin{table}[H]
    \centering
    \small
    \caption{Resultados del experimento A/B: costo, latencia y registros}
    \label{tab:sprint_ab_costo}
    \begin{tabular}{|p{3.3cm}|p{3.2cm}|p{4.6cm}|p{3.4cm}|}
        \hline
        \textbf{ID} & \textbf{Qué cambió} & \textbf{Costo / latencia} & \textbf{Log} \\
        \hline
        \texttt{UNet\_\allowbreak{}s1f11\_\allowbreak{}raw} & nada & 5.8 s/época; 336--648 s por corrida; inferencia sin costo extra & \texttt{outputs/\allowbreak{}s1f11\_\allowbreak{}src\_\allowbreak{}unet/\allowbreak{}} \\
        \hline
        \texttt{UNet\_\allowbreak{}s1f11\_\allowbreak{}trHM} & histogramas igualados a Perú & $+0.11$ s/recorte en CPU ($\approx 25$ s/época con 4 procesos de carga); inferencia igual & \texttt{outputs/\allowbreak{}s1f11\_\allowbreak{}trhm\_\allowbreak{}unet/\allowbreak{}} \\
        \hline
        \texttt{UNet\_\allowbreak{}s1f11\_\allowbreak{}trFDA\_\allowbreak{}b0.01} & amplitud de Fourier de Perú & $+0.26$ s/recorte en CPU; inferencia igual & \texttt{outputs/\allowbreak{}s1f11\_\allowbreak{}trfda\_\allowbreak{}unet/\allowbreak{}} \\
        \hline
    \end{tabular}
\end{table}

El costo de GPU es de 0.89~USD/h (RTX 4090). Las fuentes son \texttt{docs/\allowbreak{}EXPERIMENT\_\allowbreak{}LOG.md}
(entradas del 2026-10-13/14 al 2026-10-16), \texttt{docs/\allowbreak{}METHODS\_\allowbreak{}NOTES.md} (2026-10-16) y
\texttt{docs/\allowbreak{}phaseD\_\allowbreak{}test.md}.

\subsection{Gráfico}

\begin{figure}[H]
    \centering
    \includegraphics[width=\textwidth]{E_IMAGENES/6_Sprint/fig_ab_datos.pdf}
    \caption{Izquierda: IoU de agua en \texttt{peru-\allowbreak{}dev} por semilla (puntos) y media $\pm$
    desviación estándar (barra), con la cota superior de la U-Net entrenada con etiquetas de Perú
    (línea discontinua). Derecha: precisión y recall de agua en dev.}
    \label{fig:sprint_ab_datos}
\end{figure}

Para clasificación se suele mostrar una curva de precisión y recall. Construirla requiere los logits
de cada modelo, que no se versionan, así que se muestran la precisión y el recall en el punto de
operación del modelo (probabilidad 0.5).

\subsection{Confirmación en prueba (evaluación única)}

\begin{table}[H]
    \centering
    \small
    \caption{Diferencias pareadas por recorte en las particiones de prueba}
    \label{tab:sprint_ab_prueba}
    \begin{tabular}{|p{8.5cm}|c|c|}
        \hline
        \textbf{Comparación} & \textbf{Diferencia} & \textbf{IC 95~\%} \\
        \hline
        Var1 $-$ \textit{baseline}, \texttt{peru-\allowbreak{}test} & $+0.222$ & [0.132, 0.309] \\
        \hline
        Var1 $-$ U-Net con etiquetas de Perú, \texttt{peru-\allowbreak{}test} & $-0.030$ & [$-0.048$, $-0.013$] \\
        \hline
        Var1 $-$ \textit{baseline}, \texttt{peru-\allowbreak{}2017-\allowbreak{}test} (101 recortes) & $-0.064$ & [$-0.083$, $-0.046$] \\
        \hline
    \end{tabular}
\end{table}

\section{Validación}

\begin{itemize}
    \item \textbf{Partición agrupada, sin fuga.} \texttt{peru-\allowbreak{}dev} y \texttt{peru-\allowbreak{}test} salen de
    bloques espaciales con \textit{buffer}: es una partición por grupos, en la que una ubicación
    nunca aparece en dos particiones. Sen1Floods11 usa sus particiones oficiales.
    \item \textbf{Divide primero, transforma después.} Los cuantiles de la igualación de histogramas y
    las amplitudes de FDA se calculan \textbf{solo} con imágenes de \texttt{peru-\allowbreak{}train}, sin leer sus
    etiquetas. Nunca se usan imágenes de \texttt{peru-\allowbreak{}dev} ni de prueba para ajustar la
    transformación.
    \item \textbf{Semillas y particiones fijas.} Las tres filas usan las semillas 0, 1 y 2 y las
    mismas particiones.
    \item \textbf{La selección no usa dev.} La mejor época se elige en Sen1Floods11 \texttt{valid};
    \texttt{peru-\allowbreak{}dev} solo compara variantes.
    \item \textbf{Prueba una sola vez.} \texttt{peru-\allowbreak{}test} y \texttt{peru-\allowbreak{}2017-\allowbreak{}test} se evaluaron el
    2026-10-16, desde configuraciones congeladas con el hash de cada checkpoint
    (\texttt{configs/\allowbreak{}frozen.yaml}).
    \item \textbf{Sin validación cruzada K-fold.} Las reglas del proyecto prohíben volver a partir
    Sen1Floods11, y un K-fold sobre bloques de Perú obligaría a reentrenar cada fila K veces en GPU
    alquilada. En su lugar se reportan 3 semillas e intervalos \textit{bootstrap} por recortes.
    \item \textbf{Ruido entre corridas.} Repetir la semilla 0 del \textit{baseline} con la misma
    configuración dio una diferencia de 0.017 de IoU en \texttt{valid} (no determinismo de cuDNN y
    precisión mixta). Las diferencias del experimento ($+0.075$ y $-0.221$ en dev) son mucho mayores.
    \item \textbf{Diferencia operativa.} Var1 y Var2 se entrenaron con 8 procesos de carga de datos en
    lugar de 4, para que la transformación en CPU no frenara el entrenamiento. Esto cambia el orden
    aleatorio del aumento de datos, no los datos ni el modelo.
\end{itemize}

\section{Conclusión y decisión}

La regla de elección entre HM y FDA se fijó \textbf{antes} de ver resultados: el IoU en dev de la
semilla 0. HM obtuvo 0.9148 y FDA 0.7251. Solo la variante elegida pasó a la evaluación única en
prueba.

\textbf{Se adopta Var1 (\texttt{UNet\_\allowbreak{}s1f11\_\allowbreak{}trHM}).}
\begin{itemize}
    \item \textbf{Mejora consistente.} Supera al \textit{baseline} en todas las semillas ($+0.072$,
    $+0.055$ y $+0.100$ en dev) y reduce la variación entre semillas (desviación estándar 0.003 frente
    a 0.022).
    \item \textbf{Corrige el error diagnosticado en las semanas 4 y 5.} La precisión sube de 0.848 a
    0.962: se elimina la mayoría de los falsos positivos del desierto. El recall baja de 0.985 a 0.946.
    \item \textbf{Se confirma en prueba.} $+0.222$ [0.132, 0.309] sobre el \textit{baseline}. Queda a
    0.030 de la U-Net entrenada con etiquetas peruanas, sin usar ninguna.
    \item \textbf{Costo viable.} La transformación cuesta 0.11~s por recorte en CPU, solo durante el
    entrenamiento; la inferencia no cambia.
\end{itemize}

\textbf{Se descarta Var2 (\texttt{UNet\_\allowbreak{}s1f11\_\allowbreak{}trFDA\_\allowbreak{}b0.01}).} Baja el IoU a 0.615 $\pm$ 0.112, es
inestable entre semillas y es la variante más costosa (0.26~s por recorte). No se evaluó en prueba,
como establece el protocolo.

\section{Reproducibilidad}

Comandos de entrenamiento (RunPod, desde la raíz del repositorio):

{\footnotesize
\begin{verbatim}
# Baseline: UNet_s1f11_raw (commit 8a8fee4)
DATASET=s1f11 RUN_NAME=s1f11_src_unet bash scripts/cloud/run_queue.sh unet "0 1 2" \
  --optimizer.init_args.lr=3e-4 --trainer.max_epochs=200

# Var1: UNet_s1f11_trHM (commit e8001a4)
DATASET=s1f11 RUN_NAME=s1f11_trhm_unet bash scripts/cloud/run_queue.sh unet "0 1 2" \
  --optimizer.init_args.lr=3e-4 --trainer.max_epochs=200 -c configs/phaseD/train_hm.yaml \
  --data.init_args.num_workers=8

# Var2: UNet_s1f11_trFDA_b0.01 (commit e8001a4)
DATASET=s1f11 RUN_NAME=s1f11_trfda_unet bash scripts/cloud/run_queue.sh unet "0 1 2" \
  --optimizer.init_args.lr=3e-4 --trainer.max_epochs=200 -c configs/phaseD/train_fda.yaml \
  --data.init_args.num_workers=8
\end{verbatim}
}

Evaluación en \texttt{peru-\allowbreak{}dev} y construcción del cuaderno:

{\footnotesize
\begin{verbatim}
bash scripts/phaseD2_dev_logits.sh                 # logits en peru-dev
uv run python scripts/phaseD_dev_table.py          # tabla de dev
uv run python notebooks/build_baseline_nb.py       # cuaderno, tablas y figuras
\end{verbatim}
}

\begin{itemize}
    \item \textbf{Logs.} Cada corrida guarda \texttt{config.yaml}, \texttt{metrics.csv} y
    \texttt{run.json}, con su commit y su estado \texttt{dirty}, en \texttt{outputs/\allowbreak{}s1f11\_\allowbreak{}src\_\allowbreak{}unet/\allowbreak{}},
    \texttt{outputs/\allowbreak{}s1f11\_\allowbreak{}trhm\_\allowbreak{}unet/\allowbreak{}} y \texttt{outputs/\allowbreak{}s1f11\_\allowbreak{}trfda\_\allowbreak{}unet/\allowbreak{}}. El registro de todas
    las corridas está en \texttt{docs/\allowbreak{}EXPERIMENT\_\allowbreak{}LOG.md}.
    \item \textbf{Instantánea de los datos.} Perú v1, congelado el 2026-10-14: 2\,990 archivos
    fijados por SHA-256 en \texttt{docs/\allowbreak{}peru\_\allowbreak{}v1\_\allowbreak{}checksums.txt}, que ocupan 5.3~GB; Sen1Floods11
    ocupa 977~MB. El 2026-10-23 los datos se reconstruyeron desde cero en una segunda máquina y
    coincidieron bit a bit.
    \item \textbf{Versiones} (\texttt{uv.lock}): Python 3.12, PyTorch 2.14.0, Lightning 2.6.6,
    TerraTorch 1.2.13 y TorchGeo 0.9.0.
\end{itemize}

\section{Amenazas a la validez}

\begin{table}[H]
    \centering
    \small
    \caption{Amenazas a la validez del experimento A/B}
    \label{tab:sprint_amenazas}
    \begin{tabular}{|p{2.6cm}|p{5.8cm}|p{6.0cm}|}
        \hline
        \textbf{Tipo} & \textbf{Amenaza} & \textbf{Mitigación o estado} \\
        \hline
        Interna & Etiquetas de UNOSAT preliminares, sin validación en campo & Sin los recortes dudosos, cada fila de prueba sube entre 0.005 y 0.016; el único cambio de orden se da entre dos filas de TerraMind, ajenas a este experimento \\
        \hline
        Interna & Var1 y Var2 usaron 8 procesos de carga y el \textit{baseline} 4 & Solo afecta el orden aleatorio del aumento de datos; la diferencia observada supera mucho el ruido entre corridas (0.017) \\
        \hline
        Externa & Un solo evento de entrenamiento en Perú & Prueba en otro evento (2017): allí Var1 queda por debajo del \textit{baseline} \\
        \hline
        De constructo & Las etiquetas de 2017 solo definen suelo seco donde MNDWI $\leq 0$, así que miden sobre todo el recall del agua clara & Se informa por separado; la decisión se limita al evento de 2023 \\
        \hline
        De conclusión estadística & 3 semillas y una sola partición espacial; sin K-fold & Diferencias pareadas con IC \textit{bootstrap} en prueba \\
        \hline
    \end{tabular}
\end{table}

\section{Riesgos y próximos pasos}

\begin{itemize}
    \item \textbf{Riesgo: deriva hacia otros eventos.} Var1 sacrifica agua tenue para eliminar
    falsos positivos. En el evento de 2017 eso la deja 0.064 por debajo del \textit{baseline}.
    \item \textbf{Acción inmediata del siguiente sprint.} Medir en \texttt{peru-\allowbreak{}dev} cuánta agua tenue
    (MNDWI entre 0 y 0.3) pierde Var1 por estrato.
\end{itemize}

\section{Guion de la demo}

\begin{table}[H]
    \centering
    \small
    \caption{Guion de la exposición de la semana 6}
    \label{tab:sprint_demo}
    \begin{tabular}{|p{3.0cm}|p{7.6cm}|p{3.8cm}|}
        \hline
        \textbf{Paso} & \textbf{Qué decir} & \textbf{Qué mostrar} \\
        \hline
        1. Contexto (1 min) & Inundación nueva en Perú sin etiquetas; solo hay Sen1Floods11 e imágenes peruanas sin etiquetar. Métrica: IoU de agua en \texttt{peru-\allowbreak{}dev}. & Cuaderno, sección inicial \\
        \hline
        2. Lo que dijeron los datos (1 min) & El EDA y la semana 5 mostraron que el problema es el desierto brillante, no las variables. & Secciones de EDA e ingeniería de características \\
        \hline
        3. Línea base (30 s) & U-Net entrenada en Sen1Floods11: 0.836, con muchos falsos positivos (precisión 0.848). & Tabla~\ref{tab:sprint_ab_metricas}, fila \textit{baseline} \\
        \hline
        4. Variantes (1 min) & Mismo modelo y configuración; solo cambian las imágenes de entrenamiento: igualar histogramas o transferir la amplitud de Fourier, con imágenes de Perú sin etiquetas. & Tabla~\ref{tab:sprint_variantes} \\
        \hline
        5. Resultado (1 min) & HM sube a 0.912 en todas las semillas, gracias a la precisión; FDA cae a 0.615 y es inestable. & Figura~\ref{fig:sprint_ab_datos} \\
        \hline
        6. Decisión (1 min) & Se adopta HM: $+0.222$ sobre el \textit{baseline} en prueba y a 0.030 de entrenar con etiquetas de Perú, sin costo en la inferencia. FDA se descarta y no se evalúa en prueba. & Tabla~\ref{tab:sprint_ab_prueba} \\
        \hline
        7. Riesgo y siguiente paso (30 s) & En otro evento (2017), HM pierde agua tenue; el próximo sprint lo mide por estrato. & Sección de riesgos \\
        \hline
    \end{tabular}
\end{table}

\section{Anexo: comparación de modelos con etiquetas de Perú}

Este anexo no forma parte del experimento A/B de la semana 6: cambia el modelo, no los datos. Los
tres modelos se entrenaron con \texttt{peru-\allowbreak{}train} (etiquetas de Perú), con la configuración oficial
de cada uno, 3 semillas y selección en \texttt{peru-\allowbreak{}dev}.

\begin{table}[H]
    \centering
    \small
    \caption{Modelos entrenados con etiquetas de Perú}
    \label{tab:sprint_anexo_modelos}
    \begin{tabular}{|p{3.6cm}|p{2.5cm}|p{2.8cm}|p{2.6cm}|p{2.2cm}|}
        \hline
        \textbf{Modelo} & \textbf{IoU \texttt{peru-\allowbreak{}dev} (media $\pm$ de)} & \textbf{IoU \texttt{peru-\allowbreak{}test} [IC 95~\%]} & \textbf{Precisión / recall \texttt{peru-\allowbreak{}test}} & \textbf{IoU \texttt{peru-\allowbreak{}2017-\allowbreak{}test} (101)} \\
        \hline
        U-Net (ResNet-34) & 0.9415 $\pm$ 0.0019 & 0.828 [0.732, 0.895] & 0.966 / 0.853 & 0.878 \\
        \hline
        TerraMind v1-B \cite{jakubik_terramind_2025} & 0.9265 $\pm$ 0.0025 & 0.815 [0.718, 0.885] & 0.960 / 0.843 & 0.892 \\
        \hline
        Prithvi-EO-2.0 300M \cite{szwarcman_prithvi-eo-2_2024} & 0.9218 $\pm$ 0.0007 & 0.824 [0.731, 0.890] & 0.966 / 0.850 & 0.878 \\
        \hline
    \end{tabular}
\end{table}

\begin{figure}[H]
    \centering
    \includegraphics[width=0.8\textwidth]{E_IMAGENES/6_Sprint/fig_anexo_modelos.pdf}
    \caption{IoU de agua en \texttt{peru-\allowbreak{}dev} de los modelos entrenados con etiquetas de Perú, por
    semilla y media $\pm$ desviación estándar.}
    \label{fig:sprint_anexo_modelos}
\end{figure}

\begin{itemize}
    \item En dev, la U-Net supera a los dos modelos fundacionales. El 73~\% y el 86~\% de esa ventaja
    está en la franja de 1~px del borde de la etiqueta.
    \item En \texttt{peru-\allowbreak{}test} las tres medias quedan a menos de 0.014 entre sí, con intervalos de
    confianza solapados. No se calcularon diferencias pareadas entre ellas, así que no se afirma un
    orden.
    \item Con etiquetas del mismo evento el modelo importa poco: un umbral de MNDWI ajustado en
    \texttt{peru-\allowbreak{}train} ya llega a 0.924 en dev (semana 5). Por eso esta semana el experimento se
    centra en los datos.
\end{itemize}
